\chapter{Randomness, counting, and the polynomial hierarchy}
\label{chap:randomness}

This chapter covers probabilistic polynomial time, the polynomial hierarchy,
and counting classes. A probabilistic machine is a nondeterministic machine
whose choice bits are uniformly random and which is observed at a fixed clock;
$\cc{BPP}$, $\cc{RP}$, $\cc{coRP}$, $\cc{ZPP}$, and $\cc{PP}$ are defined from
its acceptance probability, and $\#\cc{P}$ from accepting leaves of the
computation tree, a different count when paths halt early. Formalized results
include the class definitions and elementary containments, the
Sipser--Lautemann theorem, $\cc{PH} \subseteq \cc{PSPACE}$,
$\cc{PP} \subseteq \cc{PSPACE}$, and Stockmeyer-style approximate counting with
pairwise-independent hashing as a finite statement. The main open directions
are the oracle characterization of the hierarchy, the closure properties of
$\cc{PP}$ through $\cc{GapP}$, and Toda's theorem.

\section{Randomized classes}

The acceptance probability $\mathrm{acceptProb}_N(x, T)$ of a machine $N$ is
the fraction of the $2^T$ choice sequences of length $T$ along which $N$ halts
with output $1$ within $T$ steps (\texttt{NTM.acceptProb}, defined in the
chapter on machine models). A path that halts after $t < T$ steps is counted
once for each of its $2^{T - t}$ extensions. The finite probability toolkit
used below is developed in the chapter on encodings: event probability, the
bound $2^{-k}$ on majority failure over $12k + 1$ independent blocks, the
invariance of acceptance probability under enlarging a clock by which every
path has halted, and the repetition machine \texttt{NTM.repeatAtTime} that
lifts majority amplification to machines. Adleman's theorem
$\cc{BPP} \subseteq \cc{P/poly}$ (\texttt{Complexity.BPP\_subset\_PPoly}) is in
the chapter on circuits.

\begin{definition}[$\cc{BPP}$]
  \label{def:BPP}
  \lean{Complexity.BPP, Complexity.BPTIME}
  \leanok
  \uses{def:ntm, def:models-accept-prob}
  $\cc{BPTIME}(T)$ is the set of languages $L$ for which there are a machine
  $N$ and $f = O(T)$ such that every path halts within $f(|x|)$ steps,
  $\mathrm{acceptProb}_N(x, f(|x|)) \ge 2/3$ for $x \in L$, and
  $\mathrm{acceptProb}_N(x, f(|x|)) \le 1/3$ for $x \notin L$.
  $\cc{BPP} = \bigcup_k \cc{BPTIME}(n^k)$.
\end{definition}

\begin{definition}[$\cc{RP}$ and $\cc{coRP}$]
  \label{def:RP}
  \lean{Complexity.RP, Complexity.RTIME, Complexity.coRP}
  \leanok
  \uses{def:ntm, def:models-accept-prob}
  $\cc{RTIME}(T)$ is defined like $\cc{BPTIME}(T)$ with acceptance probability
  at least $1/2$ on $x \in L$ and equal to $0$ on $x \notin L$.
  $\cc{RP} = \bigcup_k \cc{RTIME}(n^k)$ and
  $\cc{coRP} = \{L : \overline{L} \in \cc{RP}\}$.
\end{definition}

\begin{definition}[$\cc{ZPP}$]
  \label{def:ZPP}
  \lean{Complexity.ZPP}
  \leanok
  \uses{def:RP, def:models-accept-prob}
  $\cc{ZPP} = \cc{RP} \cap \cc{coRP}$. The Las Vegas characterization is a
  separate planned node (Theorem~\ref{thm:randomness-zpp-las-vegas}); the
  library has only fixed-clock probabilistic semantics.
\end{definition}

\begin{definition}[$\cc{PP}$]
  \label{def:PP}
  \lean{Complexity.PP, Complexity.PPTIME}
  \leanok
  \uses{def:ntm, def:models-accept-prob}
  $\cc{PPTIME}(T)$ is the set of languages $L$ for which there are a machine
  $N$ and $f = O(T)$ such that every path halts within $f(|x|)$ steps and
  $x \in L \iff \mathrm{acceptProb}_N(x, f(|x|)) > 1/2$.
  $\cc{PP} = \bigcup_k \cc{PPTIME}(n^k)$.
\end{definition}

\begin{theorem}[Elementary containments]
  \label{thm:randomness-containments}
  \lean{Complexity.P_subset_BPP, Complexity.BPP_subset_PP, Complexity.RP_subset_NP, Complexity.ZPP_subset_RP, Complexity.ZPP_subset_coRP, Complexity.ZPP_subset_NP}
  \leanok
  \uses{def:P, def:NP, def:BPP, def:RP, def:ZPP, def:PP}
  $\cc{P} \subseteq \cc{BPP} \subseteq \cc{PP}$; $\cc{RP} \subseteq \cc{NP}$;
  $\cc{ZPP} \subseteq \cc{RP}$, $\cc{ZPP} \subseteq \cc{coRP}$, and
  $\cc{ZPP} \subseteq \cc{NP}$.
\end{theorem}
\begin{proof}
  \leanok
  A deterministic machine viewed as probabilistic accepts with probability
  $0$ or $1$. Since $2/3 > 1/2 > 1/3$, bounded error implies unbounded error.
  A positive acceptance probability exhibits an accepting path, and
  probability $0$ excludes one, so $\cc{RTIME}(T) \subseteq \cc{NTIME}(T)$.
\end{proof}

\begin{theorem}[$\cc{P} \subseteq \cc{ZPP}$]
  \label{thm:randomness-P-subset-ZPP}
  \uses{def:P, def:ZPP}
  $\cc{P} \subseteq \cc{RP} \cap \cc{coRP} = \cc{ZPP}$.
\end{theorem}
\begin{proof}
  \uses{prop:polytime-P-boolean, thm:randomness-containments}
  A deterministic decider accepts yes-instances with probability $1$ and
  no-instances with probability $0$, so $\cc{P} \subseteq \cc{RP}$; closure of
  $\cc{P}$ under complement gives $\cc{P} \subseteq \cc{coRP}$.
\end{proof}

\begin{theorem}[$\cc{BPP}$ is closed under complement]
  \label{thm:randomness-bpp-compl}
  \uses{def:BPP}
  If $L \in \cc{BPP}$ then $\overline{L} \in \cc{BPP}$.
\end{theorem}
\begin{proof}
  Flip the output bit of the probabilistic machine, which exchanges acceptance
  probability $p$ with $1 - p$ at the same clock. (The library's proof of
  Sipser--Lautemann does not need this; it proves the $\Pi^p_2$ half
  directly.)
\end{proof}

\begin{theorem}[Class-level error reduction]
  \label{thm:randomness-bpp-amplification}
  \uses{def:BPP}
  For every $L \in \cc{BPP}$ and every polynomial $q$ there is a
  polynomial-time probabilistic machine, all of whose paths halt within its
  clock, that accepts every $x \in L$ with probability at least
  $1 - 2^{-q(|x|)}$ and every $x \notin L$ with probability at most
  $2^{-q(|x|)}$.
\end{theorem}
\begin{proof}
  \uses{thm:encodings-repeat-amplification}
  The repetition machine \texttt{NTM.repeatAtTime} with $12s + 1$ runs
  already reaches error $2^{-s}$
  (\texttt{NTM.repeatAtTime\_acceptProb\_ge\_one\_sub\_two\_pow},
  \texttt{NTM.repeatAtTime\_acceptProb\_le\_two\_pow}), but it fixes the
  repetition count and the clock in advance. The remaining work is a variant
  that computes $q(|x|)$ and the clock from the input before repeating, or a
  padding argument that reduces to a fixed count per input length.
\end{proof}

\begin{theorem}[Direct characterization of $\cc{coRP}$]
  \label{thm:randomness-corp-direct}
  \uses{def:RP}
  $L \in \cc{coRP}$ if and only if some polynomial-time probabilistic machine,
  all of whose paths halt within its clock, accepts every $x \in L$ with
  probability $1$ and every $x \notin L$ with probability at most $1/2$.
\end{theorem}
\begin{proof}
  Flip the output bit of an $\cc{RP}$ machine for $\overline{L}$, and
  conversely.
\end{proof}

\begin{theorem}[$\cc{ZPP}$ is zero-error probabilistic time]
  \label{thm:randomness-zpp-las-vegas}
  \uses{def:ZPP}
  \notready
  $\cc{ZPP}$ equals the class of languages decided by a polynomial-time
  probabilistic machine that outputs $1$, $0$, or ``don't know'', is never
  wrong, and says ``don't know'' with probability at most $1/2$; and it equals
  the class decided with zero error in expected polynomial time.
\end{theorem}
\begin{proof}
  \uses{thm:randomness-corp-direct}
  Run the $\cc{RP}$ and $\cc{coRP}$ machines and answer when one of them is
  conclusive. Neither the three-valued output convention nor an expected-time
  model is defined yet, and the expected-time semantics must be kept separate
  from the strict all-paths clock: a machine that is always correct and halts
  within a polynomial bound on every random string is already deterministic.
\end{proof}

\section{The polynomial hierarchy}

\begin{definition}[The polynomial hierarchy]
  \label{def:PH}
  \lean{Complexity.PH, Complexity.SigmaP, Complexity.PiP, Complexity.polyExistsLang, Complexity.polyForallLang, Complexity.polyExistsClass, Complexity.polyForallClass}
  \leanok
  \uses{def:P, def:pair}
  For a polynomial $p$ and a language $L'$, let
  $\exists^p L' = \{x : \exists w\ (|w| \le p(|x|) \wedge \langle x, w \rangle
  \in L')\}$ and $\forall^p L' = \{x : \forall w\ (|w| \le p(|x|) \Rightarrow
  \langle x, w \rangle \in L')\}$. For a class $\mathcal{C}$, let
  $\exists \mathcal{C}$ be the set of languages $\exists^p L'$ with $p$ a
  polynomial and $L' \in \mathcal{C}$, and similarly $\forall \mathcal{C}$.
  Then
  $\Sigma^p_0 = \cc{P}$, $\Pi^p_k = \{L : \overline{L} \in \Sigma^p_k\}$,
  $\Sigma^p_{k+1} = \exists \Pi^p_k$, and
  $\cc{PH} = \bigcup_k \Sigma^p_k$.
\end{definition}

\begin{proposition}[Structure of the levels]
  \label{prop:randomness-ph-levels}
  \lean{Complexity.PiP_zero, Complexity.SigmaP_succ, Complexity.PiP_succ, Complexity.SigmaP_one, Complexity.PiP_one, Complexity.SigmaP_subset_SigmaP_succ, Complexity.PiP_subset_PiP_succ, Complexity.P_subset_PH}
  \leanok
  \uses{def:PH}
  $\Pi^p_0 = \cc{P}$, $\Pi^p_{k+1} = \forall \Sigma^p_k$,
  $\Sigma^p_1 = \exists \cc{P}$, and $\Pi^p_1 = \forall \cc{P}$. For every
  $k$, $\Sigma^p_k \subseteq \Sigma^p_{k+1}$ and
  $\Pi^p_k \subseteq \Pi^p_{k+1}$; in particular $\cc{P} \subseteq \cc{PH}$.
\end{proposition}
\begin{proof}
  \leanok
  \uses{prop:polytime-P-boolean, thm:polytime-preimage}
  The complement of a bounded existential is the bounded universal of the
  complement, and $\cc{P}$ is closed under complement. The two level
  inclusions are proved together by induction; the base case uses the zero
  witness bound and the fact that the first projection of a pair is in
  $\cc{FP}$.
\end{proof}

\begin{theorem}[Collapse of the hierarchy]
  \label{thm:randomness-ph-collapse}
  \uses{def:PH, def:P, def:NP}
  If $\Sigma^p_k = \Pi^p_k$ for some $k \ge 1$, then $\cc{PH} = \Sigma^p_k$.
  In particular, if $\cc{P} = \cc{NP}$ then $\cc{PH} = \cc{P}$.
\end{theorem}
\begin{proof}
  \uses{prop:randomness-ph-levels, thm:polytime-np-sigma1, thm:polytime-coNP-structure}
  Adjacent existential quantifiers merge into one by pairing their witnesses,
  so $\exists \Sigma^p_k = \Sigma^p_k$; then induct on the level. The special
  case uses $\cc{NP} = \Sigma^p_1$ and $\cc{P} = \cc{NP} \Rightarrow
  \cc{NP} = \cc{coNP}$.
\end{proof}

\begin{theorem}[Oracle characterization of the hierarchy]
  \label{thm:randomness-ph-oracle}
  \uses{def:PH, def:NP}
  \notready
  For every $k$, $\Sigma^p_{k+1} = \cc{NP}^{\Sigma^p_k}$, where
  $\cc{NP}^{\mathcal{C}}$ is the class of languages decided in polynomial time
  by a nondeterministic oracle machine with an oracle in $\mathcal{C}$.
\end{theorem}
\begin{proof}
  \uses{thm:polytime-np-sigma1, prop:randomness-ph-levels}
  The library has deterministic Boolean-oracle machines but no
  nondeterministic oracle machines or relativized classes yet, and the result
  depends on the exact query model. Given those, guess the query answers
  together with $\Sigma^p_k$ certificates for the positive answers and
  $\Pi^p_k$ certificates for the negative ones.
\end{proof}

\section{The Sipser--Lautemann theorem}

\begin{lemma}[Lautemann's covering lemma]
  \label{lem:randomness-covering}
  \lean{Complexity.Lautemann.Covers, Complexity.Lautemann.exists_covers_of_eventProb_compl_le, Complexity.Lautemann.not_covers_of_eventProb_le}
  \leanok
  For $E \subseteq \{0,1\}^m$, say shifts $u_1, \dots, u_t \in \{0,1\}^m$ cover
  if every $r \in \{0,1\}^m$ has some $i$ with $r \oplus u_i \in E$. If
  $\Pr[\overline{E}] \le 2^{-k}$ and $m < kt$, some $t$ shifts cover. If
  $\Pr[E] \le 2^{-k}$ and $t < 2^k$, no $t$ shifts cover.
\end{lemma}
\begin{proof}
  \leanok
  Shifting is a measure-preserving involution of the seed space; count the
  shift tuples that fail to cover a given seed, and apply a union bound.
\end{proof}

\begin{lemma}[The Lautemann matrix is polynomial-time]
  \label{lem:randomness-sl-matrix}
  \lean{Complexity.MatrixVerdictInFP, Complexity.matrixVerdictInFP}
  \leanok
  \uses{def:FP, def:ntm}
  For every machine $N$, polynomial time bound $p$, and polarity $b$, the
  verdict of the quantifier-free matrix of Lautemann's characterization is in
  $\cc{FP}$. On input $\langle \langle x, w \rangle, r \rangle$ it decodes a
  tuple of shifts from $w$ and a seed from $r$ and returns the disjunction over
  the shifts of the majority vote of the amplified trials of $N$ on $x$, run
  for $p(|x|)$ steps along the shifted seed.
\end{lemma}
\begin{proof}
  \leanok
  \uses{thm:cobham}
  No machine is constructed. The verdict is written in Cobham's algebra: payload
  scanners decode the triple, the trial count and seed length are
  $\mathrm{smash}$ lengths built from $p$, and each trial is one run of $N$
  along its block of choice bits, simulated inside the algebra.
\end{proof}

\begin{theorem}[Sipser--Lautemann]
  \label{thm:randomness-sipser-lautemann}
  \lean{Complexity.SipserLautemann, Complexity.sipserLautemann, Complexity.BPP_subset_SigmaP_two, Complexity.BPP_subset_PiP_two}
  \leanok
  \uses{def:BPP, def:PH}
  $\cc{BPP} \subseteq \Sigma^p_2 \cap \Pi^p_2$.
\end{theorem}
\begin{proof}
  \leanok
  \uses{lem:randomness-covering, lem:randomness-sl-matrix, lem:polytime-decision-fn, thm:encodings-amplification, lem:encodings-acceptprob}
  Replace the machine's time bound by a dominating polynomial, which does not
  change acceptance probabilities once every path has halted, and amplify by
  majority vote over independent blocks. The covering lemma gives
  $x \in L \iff \exists$ shifts $\forall$ seeds, some shifted amplified trial
  accepts; the complementary form characterizes $x \notin L$ and gives the
  $\Pi^p_2$ half directly. The matrix is in $\cc{P}$ by
  Lemma~\ref{lem:randomness-sl-matrix} and
  Lemma~\ref{lem:polytime-decision-fn}.
\end{proof}

\begin{corollary}[$\cc{BPP} \subseteq \cc{PH}$]
  \label{cor:randomness-bpp-ph}
  \lean{Complexity.BPP_subset_PH}
  \leanok
  \uses{def:BPP, def:PH}
  $\cc{BPP} \subseteq \cc{PH}$.
\end{corollary}
\begin{proof}
  \leanok
  \uses{thm:randomness-sipser-lautemann}
  $\Sigma^p_2 \subseteq \cc{PH}$.
\end{proof}

\section{Counting classes}

\begin{definition}[$\#\cc{P}$]
  \label{def:sharpP}
  \lean{Complexity.SharpP, Complexity.NTM.acceptLeafCount}
  \leanok
  \uses{def:ntm}
  $\mathrm{acceptLeafCount}_N(x, T)$ is the number of accepting leaves of the
  computation tree of $N$ on $x$ truncated at depth $T$; a path that halts at
  depth $t < T$ is one leaf, unlike in $\mathrm{acceptCount}$. $\#\cc{P}$ is
  the set of $f \colon \{0,1\}^* \to \mathbb{N}$ for which some machine $N$ and
  some $T = O(n^k)$ satisfy: every path halts within $T(|x|)$ steps and
  $f(x) = \mathrm{acceptLeafCount}_N(x, T(|x|))$.
\end{definition}

\begin{lemma}[Accepting leaves]
  \label{lem:randomness-leaf-count}
  \lean{Complexity.NTM.acceptLeafCount_le, Complexity.NTM.acceptLeafCount_eq_of_le_of_allPathsHaltIn, Complexity.SharpP.le_two_pow}
  \leanok
  \uses{def:sharpP}
  $\mathrm{acceptLeafCount}_N(x, T) \le 2^T$. If every path halts within
  $T(|x|)$ steps and $T \le T'$ pointwise, the counts at $T'(|x|)$ and
  $T(|x|)$ agree, so a $\#\cc{P}$ value does not depend on the choice of a
  sufficient clock. Every $f \in \#\cc{P}$ satisfies $f(x) \le 2^{T(|x|)}$ for
  some polynomially bounded $T$.
\end{lemma}
\begin{proof}
  \leanok
  Induction on the depth of the truncated tree.
\end{proof}

\begin{definition}[$\cc{GapP}$]
  \label{def:randomness-gapp}
  \lean{Complexity.GapP}
  \leanok
  \uses{def:sharpP}
  $\cc{GapP}$ is the set of $h \colon \{0,1\}^* \to \mathbb{Z}$ of the form
  $h = f - g$ with $f, g \in \#\cc{P}$.
\end{definition}

\begin{proposition}[$\cc{GapP}$ is closed under negation]
  \label{prop:randomness-gapp-neg}
  \lean{Complexity.GapP.neg_mem}
  \leanok
  \uses{def:randomness-gapp}
  If $h \in \cc{GapP}$ then $-h \in \cc{GapP}$.
\end{proposition}
\begin{proof}
  \leanok
  Swap the two $\#\cc{P}$ functions.
\end{proof}

\begin{theorem}[$\cc{PP}$ via $\cc{GapP}$]
  \label{thm:randomness-pp-gapp}
  \uses{def:PP, def:randomness-gapp}
  $\cc{GapP}$ is closed under addition and multiplication, and
  $L \in \cc{PP}$ if and only if there is $h \in \cc{GapP}$ with
  $x \in L \iff h(x) > 0$ for all $x$.
\end{theorem}
\begin{proof}
  \uses{lem:randomness-leaf-count, lem:randomness-pp-integer, prop:randomness-gapp-neg}
  $\cc{PP}$ counts fixed-length random strings while $\#\cc{P}$ counts leaves.
  Pad the machine so that every path halts at exactly the clock, after which
  the two counts agree, and take $h(x) = 2\,\mathrm{acceptCount}(x, T) - 2^T$.
  For the converse, realize the positive and negative parts of $h$ by one
  machine with balanced branches.
\end{proof}

\begin{theorem}[$\cc{PP}$ is closed under complement]
  \label{thm:randomness-pp-compl}
  \uses{def:PP}
  If $L \in \cc{PP}$ then $\overline{L} \in \cc{PP}$.
\end{theorem}
\begin{proof}
  \uses{thm:randomness-pp-gapp}
  Flipping the output bit turns ``$> 1/2$'' into ``$\ge 1/2$''; ties are
  removed by passing to the gap $1 - 2h$, which is positive exactly when
  $h \le 0$.
\end{proof}

\begin{theorem}[Beigel--Reingold--Spielman]
  \label{thm:randomness-pp-intersection}
  \uses{def:PP}
  $\cc{PP}$ is closed under intersection and union.
\end{theorem}
\begin{proof}
  \uses{thm:randomness-pp-gapp, thm:randomness-pp-compl}
  Approximate the sign function by low-degree rational functions and evaluate
  them on $\cc{GapP}$ values.
\end{proof}

\begin{theorem}[Toda]
  \label{thm:randomness-toda}
  \uses{def:PH, def:sharpP}
  \notready
  $\cc{PH} \subseteq \cc{P}^{\#\cc{P}}$.
\end{theorem}
\begin{proof}
  \uses{thm:randomness-pp-gapp, def:randomness-hash, thm:randomness-ph-oracle}
  Needs polynomial-time machines with a function oracle, the Valiant--Vazirani
  isolation lemma (from pairwise-independent hashing), modular counting, and
  polynomial interpolation.
\end{proof}

\section{Polynomial-space upper bounds}

The two containments below are proved by bespoke enumerating machines. They
could instead be routed through the polynomial-space iteration theorem
\texttt{SpaceIter.mem\_PSPACE\_of\_iterate}, as Savitch's theorem and
$\cc{IP} \subseteq \cc{PSPACE}$ are.

\begin{lemma}[$\cc{PP}$ as an integer comparison]
  \label{lem:randomness-pp-integer}
  \lean{Complexity.PP_integer_characterization}
  \leanok
  \uses{def:PP}
  For every $L \in \cc{PP}$ there are a machine $N$ and a clock $f = O(n^m)$
  such that every path halts within $f(|x|)$ steps and
  $x \in L \iff 2^{f(|x|)} < 2 \cdot \mathrm{acceptCount}_N(x, f(|x|))$.
\end{lemma}
\begin{proof}
  \leanok
  Clear the denominator in $\mathrm{acceptCount}/2^f > 1/2$.
\end{proof}

\begin{theorem}[$\cc{PP} \subseteq \cc{PSPACE}$]
  \label{thm:randomness-pp-pspace}
  \lean{Complexity.PP_subset_PSPACE}
  \leanok
  \uses{def:PP, def:PSPACE}
  $\cc{PP} \subseteq \cc{PSPACE}$.
\end{theorem}
\begin{proof}
  \leanok
  \uses{lem:randomness-pp-integer}
  A deterministic machine enumerates the $2^{p(|x|)}$ choice sequences on a
  counter tape that doubles as the choice tape, simulates one path at a time,
  updates two tallies, and wipes its scratch space between passes; the loop is
  exponentially long but each pass stays in the same polynomial window.
\end{proof}

\begin{theorem}[$\cc{PH} \subseteq \cc{PSPACE}$]
  \label{thm:randomness-ph-subset-pspace}
  \lean{Complexity.PH_subset_PSPACE}
  \leanok
  \uses{def:PH, def:PSPACE}
  $\cc{PH} \subseteq \cc{PSPACE}$.
\end{theorem}
\begin{proof}
  \leanok
  \uses{prop:randomness-ph-levels}
  By induction on the level it suffices that $\cc{PSPACE}$ is closed under
  complement and under polynomially bounded existential quantification
  (\texttt{Complexity.PH\_subset\_PSPACE\_of}). Complement reruns the machine
  and flips the verdict with one extra cell; the existential enumerates the
  bounded witnesses on one reused work tape.
\end{proof}

\begin{corollary}[$\cc{BPP} \subseteq \cc{PSPACE}$]
  \label{cor:randomness-bpp-pspace}
  \uses{def:BPP, def:PSPACE}
  $\cc{BPP} \subseteq \cc{PSPACE}$.
\end{corollary}
\begin{proof}
  \uses{thm:randomness-containments, thm:randomness-pp-pspace}
  $\cc{BPP} \subseteq \cc{PP} \subseteq \cc{PSPACE}$; this composition is not
  yet stated in the library.
\end{proof}

\section{Hashing and approximate counting}

\begin{definition}[Pairwise-independent hash family]
  \label{def:randomness-hash}
  \lean{Complexity.PairwiseIndependentHash, Complexity.PairwiseIndependentHash.affine}
  \leanok
  \uses{def:bitstring}
  A pairwise-independent family from $n$-bit to $m$-bit strings with $s$-bit
  seeds is a map $h \colon \{0,1\}^s \times \{0,1\}^n \to \{0,1\}^m$ such that,
  over a uniform seed, $\Pr[h(x) = y] = 2^{-m}$ for all $x, y$, and
  $\Pr[h(x) = y \wedge h(x') = y'] = 2^{-2m}$ for all $x \neq x'$ and all
  $y, y'$. The affine family $x \mapsto Ax + b$ over $\mathrm{GF}(2)$, with
  $s = m(n+1)$, is one.
\end{definition}

\begin{theorem}[Relative approximate counting by hashing]
  \label{thm:randomness-approx-count}
  \lean{Complexity.ApproximateCounting.IsRelativeApproximation, Complexity.ApproximateCounting.Relative.one_sub_two_pow_le_eventProb_successEvent}
  \leanok
  \uses{def:randomness-hash}
  Let $S \subseteq \{0,1\}^n$ be finite, $\varepsilon \ge 1$, and
  $t \ge 0$. With probability at least $1 - 2^{-t}$ over its seed, the
  amplified hashing estimator returns $e$ with
  $(\varepsilon - 1)|S| \le \varepsilon e \le (\varepsilon + 1)|S|$. The
  estimator queries $S$ only through exact answers to ``is this hash cell of
  $S$ nonempty?''
\end{theorem}
\begin{proof}
  \leanok
  Stockmeyer's scheme: a factor-$16$ estimate from cell-occupancy answers at
  logarithmically many hash widths, with majority over independent affine
  hashes for each answer, applied to a Cartesian power of $S$ and followed by
  an integer root.
\end{proof}

\begin{theorem}[Stockmeyer]
  \label{thm:randomness-stockmeyer}
  \uses{def:sharpP}
  \notready
  For every $f \in \#\cc{P}$ and polynomial $q$, a probabilistic
  polynomial-time machine with an $\cc{NP}$ oracle outputs, with probability
  at least $2/3$, a value $e$ with
  $(1 - 1/q(|x|)) f(x) \le e \le (1 + 1/q(|x|)) f(x)$.
\end{theorem}
\begin{proof}
  \uses{thm:randomness-approx-count, lem:polytime-guess-verify}
  Apply Theorem~\ref{thm:randomness-approx-count} to the set of accepting
  paths, answering each cell-occupancy question with one $\cc{NP}$ query.
  Probabilistic oracle machines are not defined yet.
\end{proof}
